% Importado desde: 2026-03-28_Grafo_Local_vs_Red_Causal_para_Protoespacio.tex
\chapter{Exploracion estructural: --- grafo local versus red causal discreta --- como candidatos de protoespacio}
\label{ch:grafo-local-vs-red-causal}

\section{Objetivo}
En el documento anterior fijamos una idea importante:
\begin{displayquote}
\textbf{el protoespacio debe describirse primero por su estructura local y su dinamica, no por su representacion espectral.}
\end{displayquote}

Eso nos deja ahora con una pregunta mas concreta:
\begin{displayquote}
\textbf{si queremos una formulacion primaria del protoespacio, conviene pensarla como un grafo local o como una red causal discreta?}
\end{displayquote}

Este capitulo compara ambas opciones con calma y cierra proponiendo una forma hibrida que, por ahora, parece la mas prometedora para nuestro programa.

\section{Dos candidatos naturales}
\subsection{Candidato A: grafo local}
Un primer candidato natural es un grafo
\begin{equation}
G=(V,E),
\end{equation}
donde:
\begin{enumerate}[label=\arabic*.]
    \item \(V\) es el conjunto de vertices;
    \item \(E\) codifica vecindades o acoplos elementales;
    \item en cada vertice vive una fibra local \(\cH_v\);
    \item y la dinamica usa solo la estructura de adyacencia del grafo.
\end{enumerate}

\subsection{Candidato B: red causal discreta}
El segundo candidato natural es una estructura de eventos con relacion causal:
\begin{equation}
\cC=(X,\prec),
\end{equation}
donde:
\begin{enumerate}[label=\arabic*.]
    \item \(X\) es el conjunto de eventos;
    \item \(x\prec y\) significa que \(x\) puede influir causalmente en \(y\);
    \item la estructura ya distingue pasado y futuro;
    \item la dinamica debe respetar esa orientacion.
\end{enumerate}

\section{Que capta bien un grafo local}
\subsection{Ventaja 1: la vecindad espacial}
Un grafo local describe muy bien una idea de \enquote{quien es vecino de quien}. Eso lo vuelve muy natural para:
\begin{enumerate}[label=\arabic*.]
    \item modelos tipo tight-binding;
    \item Hamiltonianos locales;
    \item walks sobre redes;
    \item construcciones donde el espacio interno vive sobre vertices.
\end{enumerate}

\subsection{Ventaja 2: la estructura algebraica local}
Es facil implementar en un grafo:
\begin{enumerate}[label=\arabic*.]
    \item matrices internas por sitio;
    \item acoplos entre vecinos;
    \item operadores de masa;
    \item reglas locales de mezcla quiral.
\end{enumerate}

Esto lo vuelve especialmente compatible con el camino que ya hicimos hacia Dirac y Weyl.

\subsection{Ventaja 3: la linealizacion efectiva}
Si el grafo es regular o casi regular, es razonable esperar:
\begin{enumerate}[label=\arabic*.]
    \item sectores colectivos bien definidos;
    \item dispersion efectiva lineal en algun regime;
    \item y una lectura mas directa del limite continuo.
\end{enumerate}

\section{Que no capta tan bien un grafo local}
\subsection{Problema 1: el tiempo no viene dado}
Un grafo de adyacencia espacial no trae por si solo una orientacion temporal. Para tener dinamica hay que agregar:
\begin{enumerate}[label=\arabic*.]
    \item un Hamiltoniano con tiempo continuo;
    \item o un operador de paso con tiempo discreto.
\end{enumerate}

En otras palabras:
\begin{displayquote}
\textbf{el grafo espacial organiza muy bien la vecindad, pero no organiza por si solo la causalidad.}
\end{displayquote}

\subsection{Problema 2: espacio y causalidad quedan separados}
Si la meta final es un proto-\emph{espacio-tiempo}, un grafo puro obliga a introducir despues la parte temporal como estructura adicional. Eso no es invalido, pero si menos ambicioso.

\section{Que capta bien una red causal}
\subsection{Ventaja 1: la causalidad es primaria}
En una red causal, la pregunta \enquote{quien puede influir a quien?} ya viene incorporada desde el comienzo. Esto es muy atractivo para nuestro proyecto, porque la causalidad fue apareciendo como una de las piezas mas profundas.

\subsection{Ventaja 2: el tiempo deja de ser externo}
Una red causal permite que el orden temporal no sea un parametro añadido desde afuera, sino parte de la estructura discreta misma.

\subsection{Ventaja 3: cercania a la idea de protoespacio-tiempo}
Si el objetivo es mas fuerte que reproducir un Hamiltoniano bonito, y apunta a una lectura realmente geometrica, una estructura causal parece conceptualmente mas cercana al blanco que un simple grafo espacial.

\section{Que no capta tan bien una red causal sola}
\subsection{Problema 1: la estructura espinorial no viene gratis}
Una red causal por si sola no dice:
\begin{enumerate}[label=\arabic*.]
    \item que fibra interna vive en cada evento;
    \item como implementar pseudospin o quiralidad;
    \item como construir matrices tipo Pauli o tipo gamma.
\end{enumerate}

\subsection{Problema 2: el paso a Dirac no es automatico}
Aunque la causalidad sea mas fundamental, el camino tecnico hacia una ecuacion de Dirac efectiva puede volverse menos directo si no hay una organizacion clara de los grados internos y de la vecindad local.

\subsection{Problema 3: la isotropia puede ser mas dificil de controlar}
En una red causal muy general, la isotropia efectiva puede ser mas delicada que en una red regular con vecindad homogénea.

\section{Comparacion directa}
\begin{center}
\begin{tabular}{@{}p{4.2cm}p{4.5cm}p{4.5cm}@{}}
\toprule
Aspecto & Grafo local & Red causal discreta \\
\midrule
Vecindad local & muy natural & natural, pero orientada \\
Espacio interno & muy facil de anexar & hay que añadirlo explicitamente \\
Causalidad primaria & no & si \\
Tiempo como parte del sustrato & no necesariamente & si, al menos como orden \\
Ruta tecnica hacia Dirac & mas directa & menos directa al principio \\
Potencial interpretativo como protoespacio-tiempo & intermedio & alto \\
\bottomrule
\end{tabular}
\end{center}

\section{La leccion importante}
La comparacion sugiere algo muy claro:
\begin{displayquote}
\textbf{ninguno de los dos candidatos, tomado aisladamente, parece suficiente para lo que queremos.}
\end{displayquote}

El grafo local es mejor para controlar la estructura algebraica y la emergencia de Dirac. La red causal es mejor para que la causalidad y el tiempo no entren como añadidos artificiales.

\section{Propuesta hibrida}
La opcion que mejor encaja con el proyecto, al menos en esta etapa, es un objeto hibrido:
\begin{equation}
\cP_{\mathrm{hyb}}=(X,\cN,\prec,\cH_{\mathrm{loc}},\cU,\cG_{\mathrm{micro}}),
\label{eq:hybrid}
\end{equation}
donde:
\begin{enumerate}[label=\arabic*.]
    \item \(X\) es el conjunto de puntos o eventos;
    \item \(\cN\) da una vecindad local tipo grafo;
    \item \(\prec\) da una orientacion causal o un orden elemental;
    \item \(\cH_{\mathrm{loc}}\) lleva la fibra interna;
    \item \(\cU\) implementa la dinamica local respetando \(\cN\) y \(\prec\);
    \item \(\cG_{\mathrm{micro}}\) recoge las simetrias exactas del objeto.
\end{enumerate}

\begin{figure}[h!]
\centering
\begin{tikzpicture}[
box/.style={draw, rounded corners, thick, align=center, minimum width=2.7cm, minimum height=0.95cm},
arr/.style={-{Latex[length=3mm]}, thick}
]
\node[box, fill=blue!8] (a) at (0,0) {\(X\)\\eventos};
\node[box, fill=green!8] (b) at (3.4,0) {\(\cN\)\\vecindad};
\node[box, fill=yellow!12] (c) at (6.8,0) {\(\prec\)\\orden causal};
\node[box, fill=orange!10] (d) at (10.2,0) {\(\cH_{\mathrm{loc}}\)\\fibra};
\node[box, fill=red!8] (e) at (13.6,0) {\(\cU\)\\dinamica};
\draw[arr] (a) -- (b);
\draw[arr] (b) -- (c);
\draw[arr] (c) -- (d);
\draw[arr] (d) -- (e);
\end{tikzpicture}
\caption{La propuesta hibrida combina vecindad tipo grafo con estructura causal elemental.}
\end{figure}

\section{Interpretacion de la propuesta}
\subsection{El grafo da la cercania}
La parte \(\cN\) sirve para decir quien interactua localmente con quien y como se propagan los grados internos.

\subsection{La relacion causal da la orientacion}
La parte \(\prec\) sirve para que el protoespacio no sea meramente espacial, sino que ya traiga una flecha elemental de influencia.

\subsection{La fibra da la posibilidad de Dirac}
La fibra local \(\cH_{\mathrm{loc}}\) es la que permite que aparezcan sectores tipo Weyl, acoplamientos de masa y eventualmente una estructura Clifford efectiva.

\section{Que gana el proyecto con esta reformulacion}
Esta reformulacion gana tres cosas muy valiosas:
\begin{enumerate}[label=\arabic*.]
    \item deja de depender conceptualmente de Brillouin;
    \item mantiene cerca la ruta tecnica hacia Dirac;
    \item y eleva la causalidad a un papel estructural mas primario.
\end{enumerate}

\section{Que sigue abierto}
Por supuesto, todavia no resolvemos varias cosas:
\begin{enumerate}[label=\arabic*.]
    \item como imponer isotropia efectiva sin recurrir enseguida a periodicidad exacta;
    \item como definir de forma sobria \(\prec\) en una construccion concreta;
    \item como reconocer cuando una red hibrida produce de verdad un solo sector Dirac robusto;
    \item y como recuperar una geometria efectiva comun para varios observables.
\end{enumerate}

\section{Conclusion}
La comparacion entre grafo local y red causal no da un ganador absoluto. Da algo mejor: una leccion estructural.

\begin{displayquote}
\textbf{Si queremos un protoespacio que siga siendo tecnicamente util para construir Dirac, pero que tambien tenga una lectura genuinamente espacio-temporal, la opcion mas prometedora no parece ser un grafo puro ni una red causal pura, sino una estructura local fibrada con vecindad y orientacion causal.}
\end{displayquote}

Esa conclusion me parece muy alineada con todo lo que fuimos encontrando:
\begin{enumerate}[label=\arabic*.]
    \item la vecindad local importa;
    \item la causalidad importa;
    \item la fibra interna importa;
    \item y la dinamica elemental no puede quedar fuera de la definicion del objeto.
\end{enumerate}

\section{Siguiente paso natural}
El siguiente paso mas natural ya no seria comparar modelos en abstracto, sino escoger una version concreta de esta estructura hibrida y preguntar:
\begin{displayquote}
\textbf{cual es el modelo minimo de red local fibrada con orientacion causal que aun puede producir un subsector Dirac isotropo?}
\end{displayquote}

Ese ya seria un paso mucho mas constructivo.
